\documentclass[fleqn]{article}

\usepackage{aima-slides}
\usepackage[utf8]{inputenc}
\usepackage[T5]{fontenc}
\usepackage{lmodern}

\begin{document}

\begin{huge}
\titleslide{Độ bất định (Uncertainty)}{Chương 14}

%%%%%%%%%%%% Slide %%%%%%%%%%%%%%%%%%%%%%%%%%%%%%%%%%%%%%%%%%%%%%%%%%%
\heading{Nội dung}

\blob Độ bất định (Uncertainty)

\blob Xác suất (Probability)

\blob Cú pháp (Syntax)

\blob Ngữ nghĩa (Semantics)

\blob Các quy tắc suy diễn (Inference rules)


%%%%%%%%%%%% Slide %%%%%%%%%%%%%%%%%%%%%%%%%%%%%%%%%%%%%%%%%%%%%%%%%%%
\heading{Độ bất định}

Cho hành động $A_t$ = rời đi sân bay trước $t$ phút chuyến bay cất cánh\\
Liệu $A_t$ có giúp tôi đến đó đúng giờ không?

Các vấn đề:\al
1) khả năng quan sát một phần (tình trạng đường xá, kế hoạch của tài xế khác, v.v.)\al
2) cảm biến bị nhiễu (Báo cáo giao thông KCBS)\al
3) sự không chắc chắn trong kết quả hành động (xịt lốp xe, v.v.)\al
4) sự phức tạp to lớn của việc mô hình hóa và dự đoán giao thông

Do đó, một cách tiếp cận logic thuần túy có thể\\
\phantom{hoặc }1) có nguy cơ sai sự thật: ``$A_{25}$ sẽ giúp tôi đến đó đúng giờ''\\
hoặc 2) dẫn đến các kết luận quá yếu để ra quyết định:\nl
``$A_{25}$ sẽ giúp tôi đến đó đúng giờ nếu không có tai nạn trên cầu\nl
và trời không mưa và lốp xe của tôi vẫn nguyên vẹn v.v. và v.v.''

(Có thể nói $A_{1440}$ sẽ giúp tôi đến đó đúng giờ một cách hợp lý\\
nhưng tôi sẽ phải ở lại qua đêm tại sân bay $\ldots$)



%%%%%%%%%%%% Slide %%%%%%%%%%%%%%%%%%%%%%%%%%%%%%%%%%%%%%%%%%%%%%%%%%%
\heading{Các phương pháp xử lý độ bất định}

Logic \u{mặc định (Default)} hoặc \u{phi đơn điệu (nonmonotonic)}:\al
  Giả sử xe của tôi không bị xịt lốp\al
  Giả sử $A_{25}$ có tác dụng trừ khi bị mâu thuẫn bởi bằng chứng\\
Vấn đề: Những giả định nào là hợp lý? Làm thế nào để xử lý sự mâu thuẫn?

\u{Các quy tắc với hệ số fudge (Fudge factors)}:\al
$A_{25} \mapsto_{0.3}$ đến đó đúng giờ\al
$Sprinkler \mapsto_{0.99} WetGrass$\al
$WetGrass \mapsto_{0.7} Rain$\\
Vấn đề: Các rắc rối với sự kết hợp, ví dụ, $Sprinkler$ gây ra $Rain$??

\u{Xác suất (Probability)}\al
  Dựa trên các bằng chứng có sẵn,\nl
    $A_{25}$ sẽ giúp tôi đến đó đúng giờ với xác suất 0.04\\
Lý thuyết cờ bạc của Mahaviracarya (thế kỷ 9), Cardamo (1565)

(Logic \u{mờ (Fuzzy)} xử lý {\em mức độ đúng đắn (degree of truth)} CHỨ KHÔNG PHẢI độ bất định, ví dụ:\al
  $WetGrass$ đúng ở mức độ 0.2)



%%%%%%%%%%%% Slide %%%%%%%%%%%%%%%%%%%%%%%%%%%%%%%%%%%%%%%%%%%%%%%%%%%
\heading{Xác suất (Probability)}

Các khẳng định xác suất {\em tóm tắt} ảnh hưởng của\al
  \u{sự lười biếng (laziness)}: thất bại trong việc liệt kê các trường hợp ngoại lệ, các điều kiện phụ, v.v.\al
  \u{sự thiếu hiểu biết (ignorance)}: thiếu các sự kiện liên quan, các điều kiện ban đầu, v.v.

Xác suất \u{chủ quan (Subjective)} hoặc \u{Bayesian}:\\
Xác suất liên hệ các mệnh đề với trạng thái tri thức của chính một người\nl
ví dụ, $P(A_{25} | \mbox{không có tai nạn nào được báo cáo}) = 0.06$

Đây \u{không} phải là các khẳng định về thế giới

Xác suất của các mệnh đề thay đổi khi có bằng chứng mới:\nl
ví dụ, $P(A_{25} | \mbox{không có tai nạn nào được báo cáo},\ \mbox{5 a.m.}) = 0.15$

(Tương tự như trạng thái kéo theo logic $KB \models \alpha$, chứ không phải là sự thật.)



%%%%%%%%%%%% Slide %%%%%%%%%%%%%%%%%%%%%%%%%%%%%%%%%%%%%%%%%%%%%%%%%%%
\heading{Đưa ra quyết định trong điều kiện bất định}

Giả sử tôi tin vào điều sau đây:
\begin{eqnarray*}
P(A_{25}\mbox{ giúp tôi đến đó đúng giờ} | \ldots) &=& 0.04 \\
P(A_{90}\mbox{ giúp tôi đến đó đúng giờ} | \ldots) &=& 0.70 \\
P(A_{120}\mbox{ giúp tôi đến đó đúng giờ} | \ldots) &=& 0.95 \\
P(A_{1440}\mbox{ giúp tôi đến đó đúng giờ} | \ldots) &=& 0.9999 
\end{eqnarray*}
Nên chọn hành động nào?

Phụ thuộc vào \u{sự ưu tiên (preferences)} của tôi đối với việc lỡ chuyến bay so với ẩm thực sân bay, v.v.

\u{Lý thuyết hữu ích (Utility theory)} được sử dụng để biểu diễn và suy diễn các sự ưu tiên

\u{Lý thuyết quyết định (Decision theory)} = lý thuyết hữu ích + lý thuyết xác suất



%%%%%%%%%%%% Slide %%%%%%%%%%%%%%%%%%%%%%%%%%%%%%%%%%%%%%%%%%%%%%%%%%%
\heading{Các tiên đề của xác suất}

Với bất kỳ mệnh đề $A$, $B$ nào

1. $0 \leq P(A) \leq 1$\\
2. $P(True) = 1$ và $P(False) = 0$\\
3. $P(A \lor B) = P(A) + P(B) - P(A\land B)$

\vspace*{0.2in}

\epsfxsize=0.45\textwidth
\fig{\file{figures}{axiom3-venn.ps}}

de Finetti (1931): một tác tử đặt cược theo các xác suất vi phạm
các tiên đề này có thể bị buộc phải đặt cược sao cho mất tiền bất kể kết quả như thế nào.




%%%%%%%%%%%% Slide %%%%%%%%%%%%%%%%%%%%%%%%%%%%%%%%%%%%%%%%%%%%%%%%%%%
\heading{Cú pháp (Syntax)}

Tương tự logic mệnh đề: các thế giới có thể được định nghĩa bằng việc gán
giá trị cho các \u{biến ngẫu nhiên (random variables)}.

Các biến ngẫu nhiên \u{mệnh đề} hoặc \u{Boolean}\al
  ví dụ: $Cavity$ (tôi có bị sâu răng không?)\\
Bao gồm các biểu thức logic mệnh đề\al
  ví dụ: $\lnot Burglary \lor Earthquake$

Biến ngẫu nhiên \u{nhiều giá trị (Multivalued)}\al
  ví dụ: $Weather$ là một trong $\<sunny,rain,cloudy,snow\>$\\
Các giá trị phải đầy đủ (exhaustive) và loại trừ lẫn nhau (mutually exclusive)

Mệnh đề được xây dựng bằng cách gán một giá trị:\al
ví dụ: $Weather \eq sunny$; hoặc $Cavity \eq true$ cho rõ ràng



%%%%%%%%%%%% Slide %%%%%%%%%%%%%%%%%%%%%%%%%%%%%%%%%%%%%%%%%%%%%%%%%%%
\heading{Cú pháp (tiếp)}

Xác suất \u{tiên nghiệm (Prior)} hoặc \u{không điều kiện (unconditional)} của các mệnh đề\al
  ví dụ: $P(Cavity) = 0.1$ và $P(Weather \eq sunny) = 0.72$\\
tương ứng với niềm tin trước khi có bất kỳ bằng chứng (mới) nào

\u{Phân phối xác suất (Probability distribution)} đưa ra các giá trị cho tất cả các phép gán có thể có:\al
  $\pv(Weather) = \<0.72,0.1,0.08,0.1\>$ (\u{đã chuẩn hóa (normalized)}, nghĩa là tổng bằng 1)

\u{Phân phối xác suất đồng thời (Joint probability distribution)} đối với một tập các biến đưa ra\\
các giá trị cho từng phép gán có thể có đối với tất cả các biến\al
  $\pv(Weather,Cavity)$ = một ma trận các giá trị $4 \times 2$:

\[\begin{array}{l|cccc}
\hfil Weather \eq & sunny & rain & cloudy & snow \\
\hline
Cavity \eq true & & & & \\
Cavity \eq false & & & &
\end{array}\]



%%%%%%%%%%%% Slide %%%%%%%%%%%%%%%%%%%%%%%%%%%%%%%%%%%%%%%%%%%%%%%%%%%
\heading{Cú pháp (tiếp)}

Xác suất \u{có điều kiện (Conditional)} hoặc \u{hậu nghiệm (posterior)}\al
  ví dụ: $P(Cavity | Toothache) = 0.8$\al
  nghĩa là, \u{\u{biết rằng $Toothache$ là tất cả những gì tôi biết}}

Ký hiệu cho phân phối có điều kiện:\al
  $\pv(Weather | Earthquake)$ = vector 2 phần tử của các vector 4 phần tử

Nếu chúng ta biết nhiều hơn, ví dụ: $Cavity$ cũng được cho trước, thì chúng ta có\al
  $P(Cavity | Toothache,Cavity) = 1$\\
Lưu ý: niềm tin ít cụ thể hơn {\em vẫn có giá trị} sau khi có thêm bằng chứng
mới đến, nhưng không phải lúc nào cũng {\em hữu ích}

Bằng chứng mới có thể không liên quan, cho phép đơn giản hóa, ví dụ:\al
$P(Cavity | Toothache,49ersWin) = P(Cavity | Toothache) = 0.8$\\
Loại suy diễn này, được phê chuẩn bởi kiến thức miền, là rất quan trọng

%%%%%%%%%%%% Slide %%%%%%%%%%%%%%%%%%%%%%%%%%%%%%%%%%%%%%%%%%%%%%%%%%%
\heading{Xác suất có điều kiện}

Định nghĩa xác suất có điều kiện:
\[
  P(A|B) = \frac{P(A\land B)}{P(B)} \mbox{ nếu } P(B) \neq 0
\]
\u{Quy tắc nhân (Product rule)} đưa ra một công thức thay thế:\al
  $P(A\land B) = P(A|B)P(B) = P(B|A)P(A)$

Một phiên bản tổng quát áp dụng cho toàn bộ phân phối, ví dụ:\al
  $\pv(Weather,Cavity) = \pv(Weather|Cavity) \pv(Cavity)$\\
(Xem như một tập hợp $4\times 2$ các phương trình, {\em không} phải nhân ma trận.)

\u{Quy tắc chuỗi (Chain rule)} được suy ra bằng cách áp dụng liên tiếp quy tắc nhân:\al
$\pv(X_1,\ldots,X_n) = \pv(X_1,\ldots,X_{n-1})\ 
                        \pv(X_n | X_1,\ldots,X_{n-1})$\nl
                    = $\pv(X_1,\ldots,X_{n-2})\ 
                        \pv(X_{n-1} | X_1,\ldots,X_{n-2})\ 
                        \pv(X_n | X_1,\ldots,X_{n-1})$\nl
                  = $\ldots$\nl
                  = $\myprod_{i\eq 1}^n \pv(X_i | X_1,\ldots,X_{i-1})$



%%%%%%%%%%%% Slide %%%%%%%%%%%%%%%%%%%%%%%%%%%%%%%%%%%%%%%%%%%%%%%%%%%
\heading{Quy tắc Bayes (Bayes' Rule)}

Quy tắc nhân $P(A\land B) = P(A|B)P(B) = P(B|A)P(A)$
\[
{}\implies \mbox{\u{Quy tắc Bayes }}  P(A|B) = \frac{P(B|A)P(A)}{P(B)}
\]
Tại sao điều này lại hữu ích???

Để đánh giá xác suất \u{chẩn đoán (diagnostic)} từ xác suất \u{nhân quả (causal)}:
\[
  P(Nguy\hat{e}n\ nh\hat{a}n|K\hat{e}t\ qu\mbox{\`{a}}) = \frac{P(K\hat{e}t\ qu\mbox{\`{a}}|Nguy\hat{e}n\ nh\hat{a}n)P(Nguy\hat{e}n\ nh\hat{a}n)}{P(K\hat{e}t\ qu\mbox{\`{a}})}
\]
Ví dụ, gọi $M$ là viêm màng não, $S$ là cứng cổ:
\[
  P(M|S) = \frac{P(S|M)P(M)}{P(S)} = \frac{0.8 \times 0.0001}{0.1} = 0.0008
\]
Lưu ý: xác suất hậu nghiệm của bệnh viêm màng não vẫn rất nhỏ!


%%%%%%%%%%%% Slide %%%%%%%%%%%%%%%%%%%%%%%%%%%%%%%%%%%%%%%%%%%%%%%%%%%
\heading{Chuẩn hóa (Normalization)}

Giả sử chúng ta muốn tính phân phối hậu nghiệm trên $A$\\
khi biết $B\eq b$, và giả sử $A$ có các giá trị có thể có $a_1 \ldots a_m$

Chúng ta có thể áp dụng quy tắc Bayes cho mỗi giá trị của $A$:\al
  $P(A\eq a_1|B\eq b) = P(B\eq b|A\eq a_1)P(A\eq a_1)/P(B\eq b)$\al
  $\ldots$\al
  $P(A\eq a_m|B\eq b) = P(B\eq b|A\eq a_m)P(A\eq a_m)/P(B\eq b)$\\
Cộng các giá trị này lại, và lưu ý rằng $\mysum_i P(A\eq a_i|B\eq b) = 1$:
\[1/P(B\eq b)  = 1/\mysum_i P(B\eq b|A\eq a_i)P(A\eq a_i)\]
Đây là \u{hệ số chuẩn hóa (normalization factor)}, hằng số theo~$i$, được ký hiệu là $\alpha$:
\[
  \pv(A|B\eq b) = \alpha \pv(B\eq b | A)\pv(A)
\]
Thông thường tính một phân phối chưa chuẩn hóa, sau đó chuẩn hóa ở cuối\al
  ví dụ: giả sử $\pv(B\eq b | A)\pv(A) = \<0.4,0.2,0.2\>$\nl
    thì $\pv(A|B\eq b) = \alpha \<0.4,0.2,0.2\> 
                        = \frac{\<0.4,0.2,0.2\>}{0.4+0.2+0.2} 
                        = \<0.5,0.25,0.25\>$

%%%%%%%%%%%% Slide %%%%%%%%%%%%%%%%%%%%%%%%%%%%%%%%%%%%%%%%%%%%%%%%%%%
\heading{Điều kiện hóa (Conditioning)}

Giới thiệu một biến làm điều kiện bổ sung:
\[
  P(X|Y) = \mysum_z P(X|Y,Z\eq z) P(Z\eq z|Y)
\]
Trực giác: thường dễ dàng hơn để đánh giá từng trường hợp cụ thể, ví dụ:\\
$P(RunOver|Cross)$\al
  = $P(RunOver|Cross,Light\eq green)P(Light\eq green|Cross)$\al
  + $P(RunOver|Cross,Light\eq yellow)P(Light\eq yellow|Cross)$\al
  + $P(RunOver|Cross,Light\eq red)P(Light\eq red|Cross)$

Khi $Y$ vắng mặt, chúng ta có \u{tổng lấy ra (summing out)} hoặc \u{lấy biên (marginalization)}:
\[
  P(X) = \mysum_z P(X|Z\eq z) P(Z\eq z) = \mysum_z P(X,Z\eq z)
\]
Nói chung, với một phân phối đồng thời trên một tập các biến, thì
phân phối trên bất kỳ tập con nào (được gọi là phân phối \u{biên (marginal)} vì
lý do lịch sử) có thể được tính bằng cách tính tổng các biến khác.


%%%%%%%%%%%% Slide %%%%%%%%%%%%%%%%%%%%%%%%%%%%%%%%%%%%%%%%%%%%%%%%%%%
\heading{Phân phối đồng thời đầy đủ (Full joint distributions)}

Một \u{mô hình xác suất hoàn chỉnh} xác định mọi mục nhập trong phân phối
đồng thời cho tất cả các biến $\mbf{X} = X_1,\ldots,X_n$\\
Nghĩa là, một xác suất cho mỗi thế giới có thể có $X_1\eq x_1,\ldots,X_n\eq x_n$

(So sánh với các lý thuyết đầy đủ trong logic.)


Ví dụ: giả sử $Toothache$ và $Cavity$ là các biến ngẫu nhiên:
\[\begin{array}{l|cc}
 & Toothache\eq true & Toothache\eq false \\
\hline
Cavity \eq true  & 0.04 & 0.06 \\
Cavity \eq false & 0.01 & 0.89
\end{array}\]
Các thế giới có thể có loại trừ lẫn nhau $\implies$ $P(w_1 \land w_2) = 0$\\
Các thế giới có thể có là đầy đủ $\implies$ $w_1 \lor \cdots \lor w_n$ là $True$\nl
do đó $\mysum_i P(w_i) = 1$



%%%%%%%%%%%% Slide %%%%%%%%%%%%%%%%%%%%%%%%%%%%%%%%%%%%%%%%%%%%%%%%%%%
\heading{Phân phối đồng thời đầy đủ (tiếp)}

1) Đối với bất kỳ mệnh đề $\phi$ nào được định nghĩa trên các biến ngẫu nhiên\nl
   $\phi(w_i)$ là đúng hoặc sai

2) $\phi$ tương đương với phép tuyển của các $w_i$ khi $\phi(w_i)$ đúng

Do đó $P(\phi) = \mysum_{\{w_i:\ \phi(w_i)\}} P(w_i)$

Nghĩa là, xác suất không điều kiện của bất kỳ mệnh đề nào đều có thể tính được
như là tổng của các mục nhập từ phân phối đồng thời đầy đủ

Xác suất có điều kiện có thể được tính theo cùng một cách như một tỷ lệ:
\[
  P(\phi|\xi) = \frac{P(\phi\land \xi)}{P(\xi)}
\]
Ví dụ: 
\[
  P(Cavity |Toothache) 
   = \frac{P(Cavity \land Toothache)}{P(Toothache)}
   = \frac{0.04}{0.04+0.01} = 0.8
\]




%%%%%%%%%%%% Slide %%%%%%%%%%%%%%%%%%%%%%%%%%%%%%%%%%%%%%%%%%%%%%%%%%%
\heading{Suy diễn từ các phân phối đồng thời}

Thông thường, chúng ta quan tâm đến \al
  phân phối đồng thời hậu nghiệm của \u{các biến truy vấn (query variables)} $\mbf{Y}$\al
  cho trước các giá trị cụ thể $\mbf{e}$ đối với \u{các biến bằng chứng (evidence variables)} $\mbf{E}$

Gọi \u{các biến ẩn (hidden variables)} là $\mbf{H} = \mbf{X} - \mbf{Y} - \mbf{E}$

Khi đó phép tính tổng các mục nhập đồng thời được yêu cầu sẽ được thực hiện bằng cách lấy tổng
các biến ẩn:
\[
\pv(\mbf{Y}|\mbf{E}\eq \mbf{e}) = \alpha \pv(\mbf{Y},\mbf{E}\eq \mbf{e})
= \alpha \mysum_{\smbf{h}} \pv(\mbf{Y},\mbf{E}\eq \mbf{e},\mbf{H}\eq \mbf{h})
\]
Các số hạng trong phép tổng là các mục nhập đồng thời vì $\mbf{Y}$, $\mbf{E}$, và $\mbf{H}$ cùng nhau vét cạn tập các biến ngẫu nhiên

Các vấn đề rõ ràng:\al
1) Độ phức tạp thời gian trong trường hợp xấu nhất là $O(d^n)$ trong đó $d$ là số ngôi (arity) lớn nhất\al
2) Độ phức tạp không gian $O(d^n)$ để lưu trữ phân phối đồng thời\al
3) Làm thế nào để tìm ra các con số cho $O(d^n)$ mục nhập???

\end{huge} 
\end{document}
